\section{Graded Summary of Catholic Teaching on Death}

The table below distinguishes the theological authority of the principal teachings and questions treated in the body. \textit{De fide} denotes truth formally revealed by God and proposed by the Church's solemn or ordinary universal Magisterium as requiring the assent of divine and Catholic faith; \textit{Sententia communis} denotes the steady consensus of the Fathers and Catholic theologians; \textit{Open theological question} denotes a matter on which revered Catholic teachers hold differing views within the bounds of orthodoxy.

\begin{threestable}{Teaching}{Grade}{Principal witnesses}
Bodily death entered the human race as the penalty of Adam's sin; had man not sinned, he would have been preserved from death by a preternatural gift of grace & \textit{De fide} & Gen 2:17; 3:19; Wis 1:13; 2:24; Rom 5:12; 6:23; Council of Trent, Sess.\ V, cans.\ 1--2 (DS 1511--1512); CCC 1008 \\
\addlinespace[0.35em]
All human beings born with original sin are subject to bodily death (with revealed exceptions of privilege and the final generation transformed at the Parousia) & \textit{De fide} (universal law) / Catholic doctrine & Heb 9:27; 1 Cor 15:51; Council of Trent, Sess.\ V, can.\ 2 (DS 1512); St.\ Augustine, \textit{City of God} XIII.1--2; St.\ Thomas, \textit{ST} I-II, q.\ 85, aa.\ 5--6 \\
\addlinespace[0.35em]
Death is the separation of the immortal rational soul from the body; the rational soul is truly, of itself, and essentially the form of the human body & \textit{De fide} & Council of Vienne (DS 902); Lateran V, \textit{Apostolici Regiminis} (DS 1440); St.\ Thomas, \textit{ST} I, q.\ 75, aa.\ 2, 6; q.\ 76, a.\ 1; CCC 1005 \\
\addlinespace[0.35em]
Because the soul is naturally ordered to the body, death is an unnatural rending of the human person, not liberation from a bodily prison; the separated soul remains incomplete until the Resurrection & Catholic doctrine / \textit{Sententia communis} & 2 Cor 5:4; St.\ Augustine, \textit{City of God} XIII.6, 16--19; St.\ Thomas, \textit{ST} I, q.\ 75, a.\ 4; q.\ 89, aa.\ 1--2; \textit{Roman Catechism} I.11; Vatican II, \textit{Gaudium et spes} 18 \\
\addlinespace[0.35em]
Christ truly died by the separation of His human soul from His body while the Divinity remained hypostatically united to both, descended to the limbo of the Fathers, and conquered death by dying & \textit{De fide} & Luke 23:46; 1 Cor 15:20--22, 54--57; Heb 2:14--15; Lateran IV, \textit{Firmiter credimus} (DS 801); St.\ Athanasius, \textit{On the Incarnation} 8--9, 20--29; St.\ Leo, \textit{Sermons} 54, 71; St.\ Thomas, \textit{ST} III, qq.\ 50, 52 \\
\addlinespace[0.35em]
Death occurs once for each human person, closes the earthly state of wayfaring and merit (\textit{status viae}), excludes reincarnation, and is followed immediately by entry into eternal retribution (heaven immediately or after purification, or immediate descent into hell for those dying in actual mortal sin) & \textit{De fide} & Heb 9:27; Eccles 11:3; Luke 16:22--26; 23:43; Benedict XII, \textit{Benedictus Deus} (DS 1000--1002); Florence, \textit{Laetentur caeli} (DS 1304--1306); St.\ John of Damascus, \textit{De fide orth.} II.4; CCC 1013, 1021--1022 \\
\addlinespace[0.35em]
Particular Judgment of each immortal soul immediately upon death & Ordinary and Universal Magisterium / \textit{Sententia fidei proxima} & Heb 9:27; Ecclus 11:28--29; Luke 16:22; \textit{Roman Catechism} I.7; CCC 1021--1022 \\
\addlinespace[0.35em]
Through Baptism the Christian dies sacramentally with Christ, turning bodily mortality from a hopeless curse into an instrument of merit and conformity to the Cross & Catholic doctrine / \textit{Sententia communis} & Rom 6:3--11; Col 2:12; St.\ Augustine, \textit{City of God} XIII.4--5; St.\ Thomas, \textit{ST} I-II, q.\ 85, a.\ 5, ad 2; III, q.\ 52, a.\ 5, ad 2; CCC 1010 \\
\addlinespace[0.35em]
Extreme Unction (Anointing of the Sick) is a true sacrament of the New Law instituted by Christ and promulgated by St.\ James, conferring grace, wiping away sins and their relics, and strengthening the soul in its final struggle & \textit{De fide} & Jas 5:14--15; Mark 6:13; Council of Trent, Sess.\ XIV, Prooemium, chs.\ 1--3, cans.\ 1--4 (DS 1694--1700, 1716--1719); \textit{Roman Catechism} II.6 \\
\addlinespace[0.35em]
Prayer and the Eucharistic Sacrifice offered by the living truly benefit the faithful departed & \textit{De fide} & 2 Mach 12:43--46; St.\ Cyril of Jerusalem, \textit{Cat.} 23.9--10; St.\ John Chrysostom, \textit{Hom.\ on Phil.} 3; St.\ Augustine, \textit{Confessions} IX.11--13; Council of Florence (\textit{Laetentur caeli}, DS 1304); Council of Trent (Sess.\ XXII, ch.\ 2 \& can.\ 3, DS 1743, 1753; Sess.\ XXV, \textit{Decree on Purgatory}, DS 1820); Roman Requiem liturgy \\
\addlinespace[0.35em]
The human body deserves reverent burial as a temple of the Holy Ghost and seed of the Resurrection, though involuntary lack of burial does not harm the soul or impede resurrection & Catholic doctrine / \textit{Sententia communis} & 1 Cor 3:16--17; 6:19--20; Tob 1:20--21; 12:12; St.\ Augustine, \textit{City of God} I.12--13; Roman Ritual burial liturgy \\
\addlinespace[0.35em]
Precise mode by which the separated soul understands natural particulars without bodily phantasms before the Resurrection & Open theological question & St.\ Thomas, \textit{ST} I, q.\ 89, aa.\ 1--5 (influx of intelligible species by divine light, yielding general and confused rather than proper and distinct natural knowledge of natural things, alongside knowledge of determinate singulars and retained intellectual habits) versus later scholastic theories \\
\end{threestable}
